\section{Analytical controls for interpreting the simulations}
\label{app:checks}
These controls specialize the supplied outcome mechanisms and fixed-feature estimators. They are explanatory calculations, not claims of new general identification or inference results. No production accuracy measurement is implied.

\subsection{Quadratic response misspecification can cancel in the contrast}
\label{app:complement}
Consider the linear exposure model with columns $(\mathbf 1,z,x,s)$ fitted to the quadratic mean
$m(z)=0.5x+z+\gamma s^2$. Since $P\mathbf 1=\mathbf 1$, complementing treatment gives $s(\mathbf 1-z)=\mathbf 1-s(z)$ and
\begin{equation}
 m(\mathbf 1-z)=m(z)+(1+\gamma)\mathbf 1-2z-2\gamma s.
\end{equation}
Suppose the endpoint contrast is estimable. Its weights satisfy
$a^\top\mathbf 1=a^\top x=0$ and $a^\top z=a^\top s=1$.
Complementing the assignment preserves the fitted column space and changes the sign of its endpoint weights: $a(\mathbf 1-z)=-a(z)$. One way to see this is that $-a(z)$ satisfies the complemented contrast equations and belongs to the same column space; the minimum-norm outcome-weight solution is unique there. It follows that
\begin{equation}
 \frac{a(z)^\top m(z)+a(\mathbf 1-z)^\top m(\mathbf 1-z)}{2}
 =1+\gamma=\tau_G.
 \label{eq:complement}
\end{equation}
Under Bernoulli($1/2$) assignment by nodes or fixed clusters, an assignment and its complement have equal probability. Thus signed conditional biases cancel over complement pairs for which the contrast is estimable. The estimable subset is itself complement invariant. This does not establish unbiasedness over nonestimable assignments under an arbitrary pseudoinverse convention. Nor does it make each realized bias zero or eliminate assignment variance. At $p\ne1/2$ the complement probabilities differ, and other nonlinear response shapes need not satisfy this identity. Reporting only average signed bias could therefore hide consequential misspecification error.

\subsection{A coupled response pair with different policy endpoints}
\label{app:coupling}
Couple the linear and hidden-tail mechanisms using the same graph, covariates, treatment assignment and noise. Let
\begin{equation}
 E=\{\max_i s_i\leq 0.85\}.
\end{equation}
On $E$, the two mechanisms generate identical observed outcomes, while their targets differ by $\gamma$. For any shared estimation rule $T$ that sees this experiment but is not told the true mechanism label, its two outputs on $E$ coincide. For any scalar $t$,
\begin{equation}
 (t-\tau_{\rm linear})^2+(t-\tau_{\rm tail})^2\geq\frac{\gamma^2}{2}.
\end{equation}
Integrating over the coupled experiment and discarding nonnegative errors outside $E$ yields
\begin{equation}
 \max\{\mathrm{MSE}_{\rm linear},\mathrm{MSE}_{\rm tail}\}
 \geq\frac{\gamma^2}{4}\Pr(E).
 \label{eq:coupling}
\end{equation}
This applies to a common estimator across the constructed pair; it does not apply to an oracle given the true response mechanism. The independent prediction-evaluation experiment is not available to $T$. Good prediction on the interior experimental exposure distribution cannot, by itself, distinguish the two endpoint responses on $E$.

For a simple $d$-regular graph under node Bernoulli($p$), each neighbor-treated count is marginally $\operatorname{Binomial}(d,p)$. A union bound gives
\begin{equation}
 \Pr(E)\geq\max\left\{0,1-n\Pr\!\left[\operatorname{Binomial}(d,p)>0.85d\right]\right\}.
\end{equation}
This uses no independence assumption between units' exposures. It is a graph-conditional bound and is not the marginal law for cluster assignment. A zero lower bound is uninformative, not evidence that $E$ is impossible.

\subsection{Exact full-exposure HT risk under a sufficient one-hop mapping}
\label{app:ht-risk}
Let $N_i$ be the set of independent assignment units in the closed neighborhood of node $i$, and $k_{ij}=|N_i\cap N_j|$. These units are nodes under individual assignment and fixed partition pieces under cluster assignment. Write $m_i^a=m_i(a\mathbf 1)$ and $Q_{ij}^{ab}=m_i^a m_j^b+\Sigma_{ij}$ for $a,b\in\{0,1\}$. Assume the one-hop exposure event reveals the corresponding endpoint mean and the additive mean-zero noise is independent of assignment, with common covariance $\Sigma$. Then HT is unbiased, so its MSE equals its variance:
\begin{align}
 \operatorname{Var}(\widehat\tau_{\rm HT})=\frac1{n^2}\sum_{i,j}\big[&Q_{ij}^{11}(p^{-k_{ij}}-1)
 +Q_{ij}^{00}\{(1-p)^{-k_{ij}}-1\}\nonumber\\
 &+(Q_{ij}^{10}+Q_{ij}^{01})\mathbf 1\{k_{ij}>0\}\big].
 \label{eq:ht-risk}
\end{align}
For same-arm exposure indicators, their joint-to-marginal probability ratio is $p^{-k_{ij}}$ or $(1-p)^{-k_{ij}}$. Opposite-arm events are incompatible when the assignment-unit sets overlap and independent otherwise. Subtracting the endpoint-product contribution to the squared target yields Equation~\eqref{eq:ht-risk}. In the full ordered double sum, the two cross terms can equivalently be written $2Q_{ij}^{10}$ by exchanging $i,j$ and using covariance symmetry. This is the form evaluated blockwise by the analysis implementation.

The formula includes outcome-noise covariance as well as randomization variation. One-hop mechanisms with the same endpoint means and covariance share this risk even when their interior responses differ. The two-hop mechanism does not meet the stipulated exposure sufficiency and is excluded from this unbiased-risk calculation. Logarithmic risk output is needed when rare-event contributions exceed ordinary floating-point range; a finite Monte Carlo sample without such events cannot certify a small HT risk.
